% front-abstract.tex – Hebrew abstract + keywords, English abstract (final round, Writer).
\chapter*{תקציר}
\addcontentsline{toc}{chapter}{תקציר}

דוח מכין זה מציג פרויקט לניהול טכנולוגיה: מעבר חצייה חכם שמטרתו להפחית חצייה באדום במעבר חצייה מרומזר.
נקודת המוצא היא מחקר תצפית שנערך בתל אביב ובחן \num{1,392} הולכי רגל. ממי שהגיעו במופע האדום חצו באדום \pct{13.5}, והולך שעמד לבדו על המדרכה חצה באדום יותר מהולך שעמד עם אחרים.
מחקרים נוספים מראים שהנטייה להפר עולה ככל שההמתנה מתארכת. הבעיה שהדוח מגדיר היא אפוא הצירוף של הולך יחיד והמתנה ארוכה.

הדוח מנתח את מאמר הדגל באופן ביקורתי. ערכי הטבלאות חושבו מחדש, והבדיקה העלתה שבע סוגיות בנתונים, ובהן יחסי סיכויים שאינם תואמים את המקדמים שפורסמו ושורה לא עקבית במודל. כיוון הממצא המרכזי נשמר, אך גודל האפקט שונה מן המודפס.

הרעיון המקורי של הפרויקט היה ספירה לאחור שמראה לממתין כמה זמן נותר לאדום. סקירת המקורות הראתה שהתקן האמריקני מתיר ספירה להולכי רגל רק במופע הפינוי, ושספירה בזמן האדום נקשרה ליותר הפרות בקרב ילדים. לכן הפרויקט שינה כיוון: חישה בשפת המדרכה סופרת ממתינים ומודדת את זמן המתנתם, והבקר מקצר את האדום בגבולות בטוחים, בעדיפות להולך יחיד; הספירה מוצגת רק במופע הפינוי, וקמפיין שליטה חברתית מציג את ההמתנה כנורמה.

הדוח מגדיר שאלת מחקר, השערות ויעדי תכן מדידים, ומפרט דרישות שלכל אחת מהן יעד ושיטת אימות, ובהן הגבולות הבטוחים לקיצור האדום ואיסור על שמירת תמונות. שלוש חלופות הושוו במטריצת החלטה משוקללת, והחלופה שנבחרה קיבלה את הציון הגבוה ביותר. עוד מוצגים חקר שוק של מוצרים קיימים, ארכיטקטורת מערכת עם פרטיות כבר בתכנון, תוכנית עבודה לשנה עם חמש חבילות עבודה, מטריצת אחריות ותרשים גאנט, וניתוח סיכונים עם הפחתה ותוכנית חלופית לכל סיכון קריטי.
הדוח מסתיים בהערכה שהרכיבים הטכניים זמינים, ושהפער העיקרי הוא אישור רגולטורי לקיצור האדום.

\paragraph{מילות מפתח:} חצייה באדום; הולכי רגל; שליטה חברתית; אדום מסתגל; ספירה לאחור במופע הפינוי; חישת הולכי רגל; פרטיות כבר בתכנון; דוח מכין.

\begin{otherlanguage}{english}
\section*{\en{Abstract}}
\en{This preparatory report proposes a smart pedestrian crossing that aims to reduce red-light crossing at signalised crosswalks.
It starts from a field study of 1,392 pedestrians in Tel Aviv, where 13.5\% of those who arrived on red crossed on red, and pedestrians waiting alone crossed more often than those waiting with others.
A critical re-analysis of the published tables finds seven data issues; the main direction of the finding holds, but the effect size differs from the printed one.
The original idea, a countdown of the remaining red time, was dropped: the US standard allows pedestrian countdowns only during the clearance interval, and a red-phase countdown was linked to more violations among children.
The chosen design detects waiting pedestrians at the curb, shortens the pedestrian red within safe limits with priority for a lone pedestrian, shows a countdown only during clearance, and adds a social-control campaign.
The report covers requirements, a weighted decision matrix, market research, architecture with privacy by design, a twelve-month work plan and a risk register.}

\en{\textbf{Keywords:} pedestrian red-light crossing; social control; adaptive signal timing; countdown signals; pedestrian detection.}
\end{otherlanguage}
